\chapter{Introduction}

\section{Motivation}

\subsection{Time-domain DP simulations}
In recent years, the offshore industry has seen a significant increase in the use 

\section{Previous work}
 - Increased use of dynamic DP simulations
 - Deterministic DP simulations, no margin for varying input parameters 
 - Time consuming simulations

 - Probabilistic approach to capability - calculated capability for different Hs/Tp/Vw combinations, then determined most probable capability based on chance of each Hs/Tp/Vw combination occuring \cite{mauroProbabilisticApproachDynamic2022}
 - Hydrodynamic load estimation using PINNs, used as added mass etc isn't perfect. Estimateed based on measured motions. \cite{dingHydrodynamicsIdentificationDynamic2025a}
 - Simulation of pipe coupled with DP \cite{aiModelingSimulationDeepwater2018}, \cite{armaogluAdvantagesUsingTimeDomain2014}, \cite{li-pingCouplingTimeDomainAnalysis2014}
 
 - Developed a fast method for evaluating the DP capability in time domain for early design. Not sure how it is fast though, this isn't discussed \cite{luebckeFastMethodEvaluation2015}

 - \cite{serrarisDPStationkeepingAccuracy} performed DP capability analyses with aNySIM from MARIN. performed 8700 simulations of each 10 minutes approximately to determine capability under various conditions. Only had 9 headings

 - Following studies quantified size of their modelling errors. All used deterministic models:
    - \cite{borhaugExperimentalValidationDynamic}, comparison to model tests. 10\% error in power demand
    - \cite{martelliTimedomainMethodologyAssess2022} performed model tests, compared static capability to dynamic. Non-linear factor differences, changed for each condition
    - \cite{donnarummaNUMERICALMODELSSHIP2015} compared high fidelity model to low fidelity, two different power trains. Much larger power peaks, peaks 25\% to 50\% power, wheras simplified had constant 20\% consumption

 These errors tended to be quite large for use in decision making.

\cite{hopPipeIntegrityPipelay2024} and \cite{devriesEfficientModellingOffshore2020} modelled pipelay dynamics. 

- System identification of non-linear systems is a large and relatively known field
    - Create model of mainly unknown system to predict behaviour.


- \cite{tsiakalosScalableProbabilisticLoad2026a} propose a pipeline to predict power grid consumption while subject to exogenous variables such as weather conditions. They use ML to make the predictions as well as quantify the uncertainty of the prediction. 
    -   Used UMAP to reduce input dimensionality, then performed clustering to identify patterns for different households
    - Train a forecaster for each cluster to allow it to specialise more.
    - uncertainty was based on the forecasters, as these output quantiles or distributions. The more certain the forecaster was, the smaller the uncertainty.
    - Main goal was to used the clustering to improve accuracy for longer forecasts compared to a global forecaster
    - Relevance to DP -> Dynamic system subject to exogenous variables. Quantifies uncertainty of prediction

So far not covered:
How DP behaviour is influenced by varying parameters, how sensitive is it to uncertainty?
Probabilistic dynamic DP perediction, quantified uncertainty (how good is the prediction likely to be?)


\section{Research question}

How can a ship's dynamic positioning behaviour efficiently be forecast while considering environmental and system uncertainty?
